\begin{lema}[Cancellation of the sixteen terms and the gap]\label{macro:brecha}
For $E=q-V$ we have
\[
 \|E\|_{\rm op}\le C_\delta(\Delta_R+a_{\rm mac}+C_{\mathcal I}/n)^{1/4}.
\]
If $\Delta$, $a_{\rm mac}$, and $C_{\mathcal I}/n$ are sufficiently small in terms of $\delta$ and
the profile, then the row--column Laplacian with weights $q$ has gap at least $c_\nu/2$ on the
orthogonal complement of the constant vector.
\end{lema}
\begin{proof}
By the singular value bound and the trace expansion,
\[
 \|E\|_{\rm op}^4\le {\rm tr}((EE^T)^2)
 =\sum_{i,i',j,j'}E_{ij}E_{i'j}E_{i'j'}E_{ij'} .
\]
From the summand we can factor out the indicator that all four row and column lines are available,
together with $k^{-4}\prod_m\nu_1(e_m)$. The remaining factor is $\prod_{m=1}^4(B_m-1)$, where

\[ B_m=\frac{X_{D(e_m)}X_{A(e_m)}}{U_{D(e_m)}U_{A(e_m)}}. \]
Writing $B_T=\prod_{m\in T}B_m$ and listing the elements of each index set as subscripts, we obtain
the full expansion
\[
 \begin{aligned}
 1
 &-(B_1+B_2+B_3+B_4)\\
 &+(B_{12}+B_{13}+B_{14}+B_{23}+B_{24}+B_{34})\\
 &-(B_{123}+B_{124}+B_{134}+B_{234})+B_{1234}.
 \end{aligned}
\]
The exceptional quadruples contribute $O_\delta(1/n)$ in total, since each product is
$O_\delta(k^{-4})$. For each ordinary term, the weight multiplying the pattern of
Definition~\ref{macro:rectangulos} in the sum normalized by $n^4$ is
\[
 W_T=x^{-4}\frac{\prod_m\nu_1(e_m)}
                 {\prod_{m\in T}U_{D(e_m)}U_{A(e_m)}} .
\]
Its supremum norm and Lipschitz constant are bounded in terms of $\delta$ and the profile. Replacing
it by a constant on each box incurs a total error of $O_\delta(a_{\rm mac})$, since for each $T$ the
masses of all boxes sum to at most one. Replacing the discrete masses by their targets then incurs
an error of $O_\delta(\Delta_R)$. Restoring the variable weight in the integral gives the same
weighted target density $\prod_m\nu_1(e_m)$ for each of the sixteen terms. Its total coefficient is
therefore $1-4+6-4+1=0$. This proves the norm estimate.

Let $V'=V-c_\nu\mathbf1\mathbf1^T/k$. Its entries are nonnegative, and its row and column sums are
at most $1-c_\nu+O_\delta(\zeta)$ by Lemma~\ref{macro:balances}. The Schur bound gives the same
upper bound for $\|V'\|_{\rm op}$. If $(u,v)\perp\mathbf1_{2k}$, then $(\sum_i u_i)(\sum_j
v_j)=-(\sum_i u_i)^2\le0$. The constant matrix thus contributes a nonnegative term to the quadratic
form of the Laplacian. The degrees on the diagonal are $1+O_\delta(\zeta)$, and
$2\|u\|_2\|v\|_2\le\|u\|_2^2+\|v\|_2^2$. The resulting lower bound for the quadratic form is
\[
 \bigl(c_\nu-O_\delta(\zeta)-\|E\|_{\rm op}\bigr)
                   (\|u\|_2^2+\|v\|_2^2).
\]
Reducing the three thresholds gives the claimed gap. The exponent $1/4$ is used only to guarantee
this gap.
\end{proof}
